% --- Back cover ---
% Large PERSONIX wordmark, no logo image. Cream paper inherited globally.
% Absolute (tikz overlay) positioning so the footer sticks to the bottom.
\clearpage
\thispagestyle{empty}
\begin{tikzpicture}[remember picture, overlay]
  % --- Wordmark ---
  \node[anchor=north, font=\sffamily\bfseries\fontsize{48}{52}\selectfont]
    at ([yshift=-26mm]current page.north) {PERSONIX};
  \node[anchor=north, font=\rmfamily\itshape\large]
    at ([yshift=-44mm]current page.north) {Ndryshim pa kompromis};
  \node[anchor=north, font=\rmfamily\small,
        text width=\paperwidth-50mm, align=center]
    at ([yshift=-54mm]current page.north)
    {Një rrjet i decentralizuar reputacioni, i pacensurueshëm dhe i pakorruptueshëm};
  % --- Body ---
  \node[anchor=north, text width=\paperwidth-50mm, align=justify, font=\small]
    at ([yshift=-72mm]current page.north)
    {%
      Shteti funksiononte kur komunikimi ishte i ngadaltë, vendimet ishin
      lokale dhe autoriteti jetonte në të njëjtin qytet që sundonte. Rrjetet e
      sotme i përmbysin të tria supozimet --- dhe institucionet e ndërtuara mbi
      to nuk i përshtaten më botës që qeverisin. Ky libër përshkruan një mjet që
      i përshtatet: një rrjet të decentralizuar reputacioni ku identiteti është
      yti për ta mbajtur, e vërteta është publikisht e auditueshme dhe ryshfeti
      është vetëvrasje reputacionale.

      \smallskip
      Jo manifest. Jo parashikim. Një projekt pune se si mund të ndërtohet
      pasardhësi i shtetit --- vullnetarisht, një shpallje pas tjetrës.
    };
  % --- Footer (anchored to the bottom of the page) ---
  \node[anchor=south, font=\sffamily\footnotesize,
        text width=\paperwidth-50mm, align=center]
    at ([yshift=12mm]current page.south)
    {%
      Pavel Kudrna\quad\textbullet\quad Pragë, 2026\par
      \href{https://personix.org}{personix.org}\par
      Licencuar sipas Creative Commons Attribution-ShareAlike 4.0 International
      (CC BY-SA 4.0).
    };
\end{tikzpicture}%
\mbox{}%
\clearpage
